\section{Utilisations de l'acide benzoïque}
L'acide benzoïque $\ce{C6H5COOH}$, connu sous le code E210, est utilisé dans de
nombreux produits pharmaceutiques et comme conservateur dans certains produits
alimentaires tel que les jus de fruits, les boissons gazeuses non alcoolisés. Il
est aussi utilisé dans la synthèse de certains esters utilisés en parfumerie.
L'acide benzoïque pur se présente sous forme de cristaux blancs. Il peut être
préparé en laboratoire selon un protocole expérimental bien déterminé.

La première partie de cet exercice vise a déterminer le pourcentage de l'acide
benzoïque pur contenu dans un échantillon préparé par un chimiste en
laboratoire, et la deuxième partie s'intéresse a la préparation d'un ester a
partir de l'acide benzoïque.


\begin{donnees}
    \begin{itemize}
        \item $\mathrm{K_{A}}(\ce{C6H5COOH_{(aq)}\ttslash{}C6H5COO^{-}_{(aq)}})=
              \num{6,31e-5}$~;
        \item $M(\ce{C6H5CO2H})=\qty{122}{\gram\per\mole}$.
    \end{itemize}
\end{donnees}

\subsection{Détermination du pourcentage d'acide benzoïque pur contenu dans un
échantillon de cristaux préparés}
Un chimiste a préparé au laboratoire une quantité de cristaux d'acide benzoïque
de masse $m_{0}=\qty{244}{\mg}$. Après l'avoir dissout totalement dans de l'eau
distillée, il a obtenu une solution aqueuse ($S_{0}$) de
volume $V_{0}=\qty{100}{\mL}$ et de $\mathrm{pH}\simeq 2,95$.
\begin{enumerate}
    \item Écrire l'équation de la réaction modélisant la transformation ayant
          lieu entre l'acide benzoïque $\ce{C6H5COOH_{(aq)}}$ et l'eau.
    \item Calculer la valeur du $\mathrm{pK_{A}}$ du couple
          $\ce{C6H5COOH_{(aq)}\ttslash{}C6H5COO^{-}_{(aq)}}$.
    \item Déterminer, en justifiant votre réponse, l'espèce du couple
          $\ce{C6H5COOH_{(aq)}\ttslash{}C6H5COO^{-}_{(aq)}}$ qui prédomine dans
          la solution $(S_{0})$.
    \item Pour connaître la valeur de la masse $m$ d'acide pur présent dans les
          cristaux préparés, le chimiste a dosé le volume
          $V_{A}=\qty{10,0}{\mL}$ de la solution ($S_{0}$) par une solution
          aqueuse d'hydroxyde de sodium $\ce{Na^{+}_{(aq)} + HO^{-}_{(aq)}}$ de
          concentration molaire $C_{B}=\qty{1,0e-2}{\M}$. Le volume ajouté à
          l'équivalence est $V_{B,E}=\qty{18,0}{\mL}$.
          \begin{enumerate}
              \item Écrire l'équation de la réaction qui se produit entre
                    l'acide benzoïque $\ce{C6H5COOH_{(aq)}}$ et les ions
                    hydroxyde $\ce{HO^{-}_{(aq)}}$ considérée comme totale.
              \item Calculer la valeur de la concentration molaire $C_{A}$ de la
                    solution ($S_{0}$) préparée.
              \item En déduire la valeur de la masse $m$ d'acide benzoïque pur
                    présent dans de la solution ($S_{0}$) de volume $V_{0}$.
              \item Déterminer la valeur du pourcentage $p$ d'acide benzoïque
                    pur contenu dans les cristaux préparés par le chimiste.
          \end{enumerate}
\end{enumerate}



\subsection{Préparation d'un ester à partir de l'acide benzoïque}
L'acide benzoïque est utilisé dans la préparation des esters odorants comme le
benzoate de méthyle $\ce{C6H5-COO-CH3}$, qui est préparé a partir de la
réaction d'estérification entre l'acide benzoïque et le méthanol en présence
d'acide sulfurique selon l'équation~:
\[
    \ce{C_6H_5-COOH + CH_3-OH <=> C_6H_5-COO-CH_3 + H_2O}
\]

On réalise l'estérification à partir d'un mélange équimolaire contenant
$n=\qty{0,3}{\mole}$ d'acide benzoïque et $n=\qty{0,3}{\mole}$ de méthanol. La
constante d'équilibre $\mathrm{K}$ associée à l'équation de la réaction
d'estérification est $\mathrm{K}=4$. 
\begin{enumerate}
    \item Citer le rôle joué par l'acide sulfurique au cours de cette réaction.
    \item Dresser le tableau d'avancement correspondant à cette réaction
          d'estérification.
    \item Montrer que l'expression de $x_{\eq}$ l'avancement de la réaction à
          l'équilibre s'écrit~:
          $x_{\eq}=\frac{n\cdot\sqrt{\mathrm{K}}}{(1+\sqrt{\mathrm{K}})}$.
    \item Déterminer la composition du mélange à l'état d'équilibre du système
          chimique.
    \item Calculer la valeur du rendement $r$ de la réaction.
    \item On ajoute une quantité d'acide benzoïque au système chimique en état
          d'équilibre.

          Répondre par Vrai ou Faux aux propositions a, b et c suivantes~:
          \begin{table}[H]
              \centering
              \renewcommand{\arraystretch}{1.5}
              \begin{tabularx}{0.8\textwidth}{|
                    >{\arraybackslash\centering}p{6mm}|L|}
                  \hline a & L'équilibre du système chimique se déplace dans le
                             sens direct \\
                  \hline b & Le rendement de cette réaction augmente \\
                  \hline c & La valeur de la constante d'équilibre $\mathrm{K}$
                             augmente \\
                  \hline
              \end{tabularx}
          \end{table}
\end{enumerate}


